% Options for packages loaded elsewhere
\PassOptionsToPackage{unicode}{hyperref}
\PassOptionsToPackage{hyphens}{url}
\PassOptionsToPackage{dvipsnames,svgnames,x11names}{xcolor}
%
\documentclass[
  10pt,
  a4paper,
]{article}
\usepackage{amsmath,amssymb}
\usepackage{iftex}
\ifPDFTeX
  \usepackage[T1]{fontenc}
  \usepackage[utf8]{inputenc}
  \usepackage{textcomp} % provide euro and other symbols
\else % if luatex or xetex
  \usepackage{unicode-math} % this also loads fontspec
  \defaultfontfeatures{Scale=MatchLowercase}
  \defaultfontfeatures[\rmfamily]{Ligatures=TeX,Scale=1}
\fi
\usepackage{lmodern}
\ifPDFTeX\else
  % xetex/luatex font selection
\fi
% Use upquote if available, for straight quotes in verbatim environments
\IfFileExists{upquote.sty}{\usepackage{upquote}}{}
\IfFileExists{microtype.sty}{% use microtype if available
  \usepackage[]{microtype}
  \UseMicrotypeSet[protrusion]{basicmath} % disable protrusion for tt fonts
}{}
\makeatletter
\@ifundefined{KOMAClassName}{% if non-KOMA class
  \IfFileExists{parskip.sty}{%
    \usepackage{parskip}
  }{% else
    \setlength{\parindent}{0pt}
    \setlength{\parskip}{6pt plus 2pt minus 1pt}}
}{% if KOMA class
  \KOMAoptions{parskip=half}}
\makeatother
\usepackage{xcolor}
\usepackage[margin=2.2cm]{geometry}
\usepackage{color}
\usepackage{fancyvrb}
\newcommand{\VerbBar}{|}
\newcommand{\VERB}{\Verb[commandchars=\\\{\}]}
\DefineVerbatimEnvironment{Highlighting}{Verbatim}{commandchars=\\\{\}}
% Add ',fontsize=\small' for more characters per line
\newenvironment{Shaded}{}{}
\newcommand{\AlertTok}[1]{\textcolor[rgb]{1.00,0.00,0.00}{\textbf{#1}}}
\newcommand{\AnnotationTok}[1]{\textcolor[rgb]{0.38,0.63,0.69}{\textbf{\textit{#1}}}}
\newcommand{\AttributeTok}[1]{\textcolor[rgb]{0.49,0.56,0.16}{#1}}
\newcommand{\BaseNTok}[1]{\textcolor[rgb]{0.25,0.63,0.44}{#1}}
\newcommand{\BuiltInTok}[1]{\textcolor[rgb]{0.00,0.50,0.00}{#1}}
\newcommand{\CharTok}[1]{\textcolor[rgb]{0.25,0.44,0.63}{#1}}
\newcommand{\CommentTok}[1]{\textcolor[rgb]{0.38,0.63,0.69}{\textit{#1}}}
\newcommand{\CommentVarTok}[1]{\textcolor[rgb]{0.38,0.63,0.69}{\textbf{\textit{#1}}}}
\newcommand{\ConstantTok}[1]{\textcolor[rgb]{0.53,0.00,0.00}{#1}}
\newcommand{\ControlFlowTok}[1]{\textcolor[rgb]{0.00,0.44,0.13}{\textbf{#1}}}
\newcommand{\DataTypeTok}[1]{\textcolor[rgb]{0.56,0.13,0.00}{#1}}
\newcommand{\DecValTok}[1]{\textcolor[rgb]{0.25,0.63,0.44}{#1}}
\newcommand{\DocumentationTok}[1]{\textcolor[rgb]{0.73,0.13,0.13}{\textit{#1}}}
\newcommand{\ErrorTok}[1]{\textcolor[rgb]{1.00,0.00,0.00}{\textbf{#1}}}
\newcommand{\ExtensionTok}[1]{#1}
\newcommand{\FloatTok}[1]{\textcolor[rgb]{0.25,0.63,0.44}{#1}}
\newcommand{\FunctionTok}[1]{\textcolor[rgb]{0.02,0.16,0.49}{#1}}
\newcommand{\ImportTok}[1]{\textcolor[rgb]{0.00,0.50,0.00}{\textbf{#1}}}
\newcommand{\InformationTok}[1]{\textcolor[rgb]{0.38,0.63,0.69}{\textbf{\textit{#1}}}}
\newcommand{\KeywordTok}[1]{\textcolor[rgb]{0.00,0.44,0.13}{\textbf{#1}}}
\newcommand{\NormalTok}[1]{#1}
\newcommand{\OperatorTok}[1]{\textcolor[rgb]{0.40,0.40,0.40}{#1}}
\newcommand{\OtherTok}[1]{\textcolor[rgb]{0.00,0.44,0.13}{#1}}
\newcommand{\PreprocessorTok}[1]{\textcolor[rgb]{0.74,0.48,0.00}{#1}}
\newcommand{\RegionMarkerTok}[1]{#1}
\newcommand{\SpecialCharTok}[1]{\textcolor[rgb]{0.25,0.44,0.63}{#1}}
\newcommand{\SpecialStringTok}[1]{\textcolor[rgb]{0.73,0.40,0.53}{#1}}
\newcommand{\StringTok}[1]{\textcolor[rgb]{0.25,0.44,0.63}{#1}}
\newcommand{\VariableTok}[1]{\textcolor[rgb]{0.10,0.09,0.49}{#1}}
\newcommand{\VerbatimStringTok}[1]{\textcolor[rgb]{0.25,0.44,0.63}{#1}}
\newcommand{\WarningTok}[1]{\textcolor[rgb]{0.38,0.63,0.69}{\textbf{\textit{#1}}}}
\usepackage{longtable,booktabs,array}
\usepackage{calc} % for calculating minipage widths
% Correct order of tables after \paragraph or \subparagraph
\usepackage{etoolbox}
\makeatletter
\patchcmd\longtable{\par}{\if@noskipsec\mbox{}\fi\par}{}{}
\makeatother
% Allow footnotes in longtable head/foot
\IfFileExists{footnotehyper.sty}{\usepackage{footnotehyper}}{\usepackage{footnote}}
\makesavenoteenv{longtable}
\setlength{\emergencystretch}{3em} % prevent overfull lines
\providecommand{\tightlist}{%
  \setlength{\itemsep}{0pt}\setlength{\parskip}{0pt}}
\setcounter{secnumdepth}{-\maxdimen} % remove section numbering
\ifLuaTeX
  \usepackage{selnolig}  % disable illegal ligatures
\fi
\usepackage{bookmark}
\IfFileExists{xurl.sty}{\usepackage{xurl}}{} % add URL line breaks if available
\urlstyle{same}
\hypersetup{
  colorlinks=true,
  linkcolor={blue},
  filecolor={Maroon},
  citecolor={Blue},
  urlcolor={blue},
  pdfcreator={LaTeX via pandoc}}

\author{}
\date{}

\begin{document}

\section{BA2-T3 - Cross-corpus regression of operator/role/modifier
candidate and semantic-key/lexical
separation}\label{ba2-t3---cross-corpus-regression-of-operatorrolemodifier-candidate-and-semantic-keylexical-separation}

\textbf{Revision:} R1

\textbf{Status:} COMPLETED / PROVISIONAL PASS WITH REDUCTION / BA2 NOT
CLOSED

\textbf{Repository baseline reviewed:}
\texttt{4d832bcf90109106d543029cb517be32a6fe7ea7}

\textbf{Phase:} BA2 - Relations and canonical action vocabulary

\textbf{BA0:} CLOSED

\textbf{BA1:} CLOSED

\textbf{BA3:} NOT STARTED

\subsection{1. Question}\label{1-question}

BA2-T3 asks only:

\begin{quote}
Can the BA2-T2 semantic-operator/role/modifier candidate replay both
governed corpora without raw-prose reconstruction, and which candidate
distinctions can be removed or promoted while preserving the semantics
needed by methodology-neutral consumers?
\end{quote}

This is a regression/reduction step. It does not reward vocabulary
growth.

\subsection{2. Evidence replayed}\label{2-evidence-replayed}

The regression uses:

\begin{itemize}
\tightlist
\item
  closed BA1 identity ontology (\texttt{BAReferent\ +\ BAProposition});
\item
  BA2-T1 n-ary proposition structure;
\item
  BA2-T2 R1 vocabulary candidate;
\item
  facial-access D-3.3, FR-3.3, FR-3.4, D-3.4, D-3.5 and SR-3.4-C/I/P;
\item
  order-fulfillment MR-1 through MR-4, including reservation, payment,
  fulfillment, provider contracts, lifecycle and physical handoff reuse;
\item
  BA0-T2 STRIDE-oriented transfer projection only as a bounded consumer
  probe.
\end{itemize}

No source document is changed by this trial.

\subsection{3. Regression strategy}\label{3-regression-strategy}

For every candidate operator the test attempts, in order:

\begin{enumerate}
\def\labelenumi{\arabic{enumi}.}
\tightlist
\item
  \textbf{delete} it and replay the two corpora;
\item
  \textbf{merge} it into another key and check for semantic loss;
\item
  \textbf{move} the distinction into roles, typed values or constraints
  where that is more economical;
\item
  retain the key only if a consumer would otherwise have to reconstruct
  relevant project meaning from prose.
\end{enumerate}

The same discipline is applied to roles, modifier kinds and
classification.

\subsection{4. Operator regression
ledger}\label{4-operator-regression-ledger}

\begin{longtable}[]{@{}lll@{}}
\toprule\noalign{}
BA2-T2 key & Regression result & Disposition \\
\midrule\noalign{}
\endhead
\bottomrule\noalign{}
\endlastfoot
\texttt{transfer} & Required for camera delivery and
order/provider/carrier conveyance; not equivalent to
reference/dependency & \textbf{RETAIN} \\
\texttt{produce} & Required to preserve governed output/result
production; \texttt{map}/\texttt{derive} can use input/result roles plus
constraints & \textbf{RETAIN} \\
\texttt{create} & Required for OrderEvaluation, Reservation,
FulfillmentTask, StockMovement/event establishment & \textbf{RETAIN} \\
\texttt{observe} & Required to distinguish reading/querying current
state from creation or mutation & \textbf{RETAIN} \\
\texttt{transition} & Required for Order/Reservation lifecycle and state
change & \textbf{RETAIN} \\
\texttt{correlate} & Same-request/evaluation identity matching is
stronger than generic reference & \textbf{RETAIN} \\
\texttt{reference} & FulfillmentTask -\textgreater{} Order/Reservation
does not imply correlation/dependency & \textbf{RETAIN} \\
\texttt{dependOn} & Required for commitment gates and physical-handoff
prerequisites & \textbf{RETAIN} \\
\texttt{consumeService} & Actual service use without ownership transfer
is not exhausted by generic dependency & \textbf{RETAIN} \\
\texttt{realize} & Connectivity -\textgreater{} Ethernet/Wi-Fi change
must remain mechanically distinct & \textbf{RETAIN} \\
\texttt{assignResponsibility} & Required for responsibility/authority
placement; can carry an authority kind & \textbf{RETAIN / BROADEN} \\
\texttt{ownOrManage} & Can be represented by
\texttt{assignResponsibility} + \texttt{responsibilityKind} + polarity
without loss & \textbf{MERGE / REMOVE KEY} \\
\texttt{constrain} & Required for security, completion, atomicity,
concurrency, idempotency and admissibility rules & \textbf{RETAIN} \\
\texttt{classify} & Required for reusable semantic kinds without BA1
subtype proliferation & \textbf{RETAIN} \\
\end{longtable}

\subsubsection{Result}\label{result}

\begin{Shaded}
\begin{Highlighting}[]
\NormalTok{BA2{-}T2 seed keys: 14}
\NormalTok{BA2{-}T3 R2 seed keys: 13}
\NormalTok{missing operator forced by regression: NONE}
\end{Highlighting}
\end{Shaded}

\subsection{\texorpdfstring{5. Why \texttt{ownOrManage} collapses but
\texttt{consumeService} does
not}{5. Why ownOrManage collapses but consumeService does not}}\label{5-why-ownormanage-collapses-but-consumeservice-does-not}

The facial corpus deliberately states both that the project
\textbf{consumes} local connectivity and that it does \textbf{not
own/manage} the underlying transport.

\texttt{consumeService} must remain distinct because it says an actual
capability/service is used. Ownership/management, however, is one
authority mode among possible responsibility modes. It can therefore be
represented under the broader responsibility relation:

\begin{Shaded}
\begin{Highlighting}[]
\NormalTok{assignResponsibility}
\NormalTok{  responsibleParty    {-}\textgreater{} project}
\NormalTok{  responsibilityScope {-}\textgreater{} underlying transport}
\NormalTok{  responsibilityKind  {-}\textgreater{} ownership/management}
\NormalTok{  polarity            {-}\textgreater{} negative}
\end{Highlighting}
\end{Shaded}

The semantic boundary
\texttt{service\ consumption\ !=\ ownership/management} remains visible
even though one operator key is removed.

\subsection{6. Missing-role test}\label{6-missing-role-test}

Two role distinctions are added because the replay cannot remain
explicit without them.

\subsubsection{\texorpdfstring{6.1 \texttt{input} under
\texttt{produce}}{6.1 input under produce}}\label{61-input-under-produce}

\texttt{PaymentAdapter\ maps\ provider\ response/outcome\ to\ PaymentAuthorizationResult}
and \texttt{AvailabilityResult} derives from current accountable
inventory state. Merely saying that a result is produced loses the
governed input/result boundary.

The regression therefore adds optional \texttt{input} participation to
\texttt{produce}.

\subsubsection{\texorpdfstring{6.2 \texttt{responsibilityKind} under
\texttt{assignResponsibility}}{6.2 responsibilityKind under assignResponsibility}}\label{62-responsibilitykind-under-assignresponsibility}

The merge of ownership/management into responsibility semantics must not
make ownership, management and generic responsibility indistinguishable.
A controlled \texttt{responsibilityKind} term preserves that
distinction.

No other missing role distinction is forced by the current corpora.

\subsection{7. Cardinality test}\label{7-cardinality-test}

The R2 role matrix uses exact maxima only where the current operator
meaning requires them. It avoids inventing general systems-model
cardinalities.

Examples:

\begin{Shaded}
\begin{Highlighting}[]
\NormalTok{transfer:}
\NormalTok{  source       1}
\NormalTok{  destination  1..*}
\NormalTok{  content      1..*}

\NormalTok{correlate:}
\NormalTok{  correlatedItem      1..*}
\NormalTok{  correlationContext  1}

\NormalTok{transition:}
\NormalTok{  subject    1}
\NormalTok{  fromState  0..1}
\NormalTok{  toState    1}
\NormalTok{  actor      0..1}

\NormalTok{consumeService:}
\NormalTok{  consumer   1..*}
\NormalTok{  service    1}
\NormalTok{  provider   0..1}
\end{Highlighting}
\end{Shaded}

The matrix is internally compatible with both corpora. No clause
requires anonymous positional arguments or method-specific DFD slots.

\subsection{8. Lexical-key separation
regression}\label{8-lexical-key-separation-regression}

The same semantic key remains stable across distinct source verbs:

\begin{itemize}
\tightlist
\item
  delivery/send/submit/emit pressure -\textgreater{} \texttt{transfer}
  where the governed meaning is conveyance;
\item
  read/query pressure -\textgreater{} \texttt{observe};
\item
  create/record pressure -\textgreater{} \texttt{create} only when new
  project-semantic identity/event is established;
\item
  transition/release/consume/update pressure -\textgreater{}
  \texttt{transition} when governed lifecycle/state changes;
\item
  produce/derive/map pressure -\textgreater{} \texttt{produce} with
  explicit inputs/results and separate constraints when normalization
  rules matter.
\end{itemize}

The reverse is also important: identical source words do not force
identical BA semantics. \texttt{record} may mean create an attributable
StockMovement or persist/update an existing state, depending on the
governed clause.

\subsubsection{Disposition}\label{disposition}

\begin{Shaded}
\begin{Highlighting}[]
\NormalTok{semantic key stable across source wording: PASS}
\NormalTok{source lexical predicate as semantic identity: REJECTED}
\NormalTok{canonical display synonym work: DEFERRED TO BA5}
\end{Highlighting}
\end{Shaded}

\subsection{9. Polarity regression}\label{9-polarity-regression}

Negative facts occur in both corpora and do not justify \texttt{notX}
operator proliferation.

Examples include:

\begin{itemize}
\tightlist
\item
  project does not own/manage transport;
\item
  incomplete transfer is not success;
\item
  stale/superseded results must not satisfy current commitment gates;
\item
  raw provider/carrier status must not directly become governed state;
\item
  rejected/indeterminate carrier acceptance is not physical handoff.
\end{itemize}

Explicit polarity therefore survives as a normalized proposition
property.

\textbf{Disposition:} \textbf{RETAIN / READY FOR BA2 CLOSURE REVIEW}.

\subsection{10. Modifier-promotion
regression}\label{10-modifier-promotion-regression}

BA2-T2 allowed a wider modifier candidate set. Replay shows that several
of those semantics are too important to remain embedded by default.

\subsubsection{\texorpdfstring{10.1 Promote to explicit
\texttt{constrain}
propositions}{10.1 Promote to explicit constrain propositions}}\label{101-promote-to-explicit-constrain-propositions}

\begin{itemize}
\tightlist
\item
  completion/failure semantics;
\item
  atomicity;
\item
  concurrency;
\item
  idempotency;
\item
  normative exact-once/all-lines quantity rules when they constrain
  behavior.
\end{itemize}

These rules recur, are independently reviewable and are useful selection
points for analysis.

\subsubsection{10.2 Keep only genuinely local
modifiers}\label{102-keep-only-genuinely-local-modifiers}

\begin{itemize}
\tightlist
\item
  \texttt{condition} (including local applicability/guard semantics);
\item
  \texttt{temporalScope}.
\end{itemize}

Participant cardinalities are structural role contracts rather than
modifiers. If a condition or temporal scope becomes independently
reusable/reviewable, the existing promotion rule moves it into
referent/proposition semantics.

\subsubsection{Disposition}\label{disposition-1}

\begin{Shaded}
\begin{Highlighting}[]
\NormalTok{untyped/free{-}text modifier bag: REJECTED}
\NormalTok{wide modifier enumeration from T2: REDUCED}
\NormalTok{condition: RETAIN}
\NormalTok{ temporalScope: RETAIN}
\NormalTok{atomicity/concurrency/idempotency/completion rules: PROMOTE TO CONSTRAIN}
\end{Highlighting}
\end{Shaded}

\subsection{11. Classification
regression}\label{11-classification-regression}

Classification-as-proposition successfully supplies reusable semantic
distinctions without turning them into BA1 root types.

A bounded consumer can identify, for example, delivery behavior,
RecognitionCapture information/resource meaning and Camera/processor
capability meaning from controlled classification assertions rather than
from source prose.

The order corpus adds pressure from party, capability,
information/resource, store and contract meanings.

However, regression does \textbf{not} justify freezing a universal
semantic-kind taxonomy. Doing so would recreate the general
systems-modeling ontology that BA1 deliberately avoided.

\subsubsection{Disposition}\label{disposition-2}

\begin{Shaded}
\begin{Highlighting}[]
\NormalTok{classification{-}as{-}proposition: RETAIN}
\NormalTok{stable semantic{-}kind registry contract: REQUIRED}
\NormalTok{fixed exhaustive semantic{-}kind enumeration: REJECTED AS CURRENT{-}SCOPE REQUIREMENT}
\NormalTok{semantic kind {-}\textgreater{} BA1 metaclass promotion: REJECTED}
\end{Highlighting}
\end{Shaded}

\subsection{12. Facial-access replay}\label{12-facial-access-replay}

The FR-3.4 branch can be selected mechanically as:

\begin{Shaded}
\begin{Highlighting}[]
\NormalTok{Delivery behavior referent}
\NormalTok{  transfer(CameraSubsystem, RecognitionProcessor, RecognitionCapture)}
\NormalTok{  correlate(RecognitionCapture, RecognitionRequest)}

\NormalTok{Connectivity meaning}
\NormalTok{  consumeService(project context, LocalConnectivity)}
\NormalTok{  assignResponsibility(project, underlying transport, ownership/management) [negative]}
\NormalTok{  realize(LocalConnectivity, WiredEthernet)}

\NormalTok{Delivery constraints}
\NormalTok{  incomplete delivery != successful completion}
\NormalTok{  confidentiality}
\NormalTok{  integrity}
\NormalTok{  authorized provenance}
\end{Highlighting}
\end{Shaded}

A future STRIDE overlay can derive source, destination, information,
correlation, external responsibility and realization without introducing
\texttt{DataFlow} or \texttt{TrustBoundary} into Base Analysis.

\subsection{13. Order-fulfillment
replay}\label{13-order-fulfillment-replay}

The order corpus does not force new operator keys:

\begin{itemize}
\tightlist
\item
  \texttt{requestsCapture}, \texttt{requestsVoid},
  \texttt{requestsStockIssue} can be represented through
  transfer/service-consumption/dependency semantics rather than copied
  as operators;
\item
  \texttt{maps} is handled by \texttt{produce} with input/result plus
  normalization constraints;
\item
  \texttt{increases}, \texttt{decreasesReserved}, \texttt{releases} and
  \texttt{consumes} are lifecycle/state transition and constraint
  semantics, not separate canonical verbs;
\item
  exactly-once, all-or-nothing, concurrency and idempotency are explicit
  constraints;
\item
  the physical handoff remains a reusable event/milestone referent and
  is not collapsed into carrier acceptance.
\end{itemize}

The corpus therefore passes without vocabulary growth.

\subsection{14. Bounded method-consumer
regression}\label{14-bounded-method-consumer-regression}

The BA0-T2 STRIDE-oriented transfer projection remains constructible
from R2:

\begin{Shaded}
\begin{Highlighting}[]
\NormalTok{behavior        \textless{}{-} classify(delivery, behavior/event)}
\NormalTok{source          \textless{}{-} transfer.source}
\NormalTok{destination     \textless{}{-} transfer.destination}
\NormalTok{information     \textless{}{-} transfer.content + classification}
\NormalTok{correlation     \textless{}{-} correlate}
\NormalTok{responsibility  \textless{}{-} consumeService + assignResponsibility/polarity}
\NormalTok{realization     \textless{}{-} realize}
\NormalTok{failure rule    \textless{}{-} constrain}
\NormalTok{security rules  \textless{}{-} constrain}
\end{Highlighting}
\end{Shaded}

STRIDE/DFD categories remain outside Base Analysis.

\textbf{Disposition:} \textbf{PASS}.

\subsection{15. BA1 reopen test}\label{15-ba1-reopen-test}

No operator, role, modifier or classification mechanism requires a third
first-class semantic identity family.

\textbf{Disposition:} \texttt{BA1\ REMAINS\ CLOSED}.

\subsection{16. Trial disposition}\label{16-trial-disposition}

\begin{Shaded}
\begin{Highlighting}[]
\NormalTok{BA2{-}T1 structural candidate                  RETAINED}
\NormalTok{exact semantic operator key                  RETAINED}
\NormalTok{operator{-}family facet                        REMOVED FROM NORMATIVE CORE}
\NormalTok{operator seed                                REDUCED 14 {-}\textgreater{} 13}
\NormalTok{ownOrManage                                  MERGED INTO assignResponsibility}
\NormalTok{missing operator                             NONE FORCED}
\NormalTok{operator{-}scoped role contracts               PASS / R2 MATRIX}
\NormalTok{new role: input                              FORCED BY REGRESSION}
\NormalTok{new role: responsibilityKind                 FORCED BY MERGE}
\NormalTok{semantic{-}key/lexical separation              PASS}
\NormalTok{explicit polarity                            RETAIN}
\NormalTok{modifier vocabulary                          REDUCED}
\NormalTok{atomicity/concurrency/idempotency/failure     PROMOTE TO constrain}
\NormalTok{classification{-}as{-}proposition                RETAIN}
\NormalTok{fixed exhaustive semantic{-}kind taxonomy      NOT REQUIRED}
\NormalTok{bounded STRIDE consumer                      PASS}
\NormalTok{third BA1 identity family                     NOT FORCED}
\NormalTok{BA2{-}T3                                        COMPLETED / PROVISIONAL PASS WITH REDUCTION}
\NormalTok{BA2                                           STARTED / NOT CLOSED}
\NormalTok{BA3                                           NOT STARTED}
\end{Highlighting}
\end{Shaded}

\subsection{17. Smallest unresolved
set}\label{17-smallest-unresolved-set}

No further discovery microstep is justified by a missing operator or
corpus counterexample. What remains is a closure review over the reduced
R2 candidate:

\begin{enumerate}
\def\labelenumi{\arabic{enumi}.}
\tightlist
\item
  accept/reject the thirteen-key current-scope seed;
\item
  accept/reject the operator-role matrix and ownership/management merge;
\item
  accept/reject polarity and the narrowed modifier boundary;
\item
  accept/reject the semantic-kind registry contract without a universal
  enumeration;
\item
  verify that the resulting BA2 contract creates no BA1 or BA3
  responsibility leakage.
\end{enumerate}

\subsection{18. Next authorized
microstep}\label{18-next-authorized-microstep}

Execute only:

\begin{quote}
\textbf{BA2-T4 - relation/action vocabulary closure review over the
regressed R2 candidate.}
\end{quote}

BA2-T4 must be a closure review, not another vocabulary-expansion
exercise. Do not start BA3, formal STRIDE overlays, Common Finding or
implementation work.

\end{document}
